\section{September 21, 2026}
% Based on pages 1--3 of 5510921.pdf.

We continue the study of implicit Runge--Kutta methods, derive their
stability functions, and develop the collocation construction outlined
on September 18. The main connection is that integrating an
interpolating polynomial produces both the RK coefficients and an
associated quadrature rule.

We retain \(x_n\), \(\tau_n=t_{n+1}-t_n\), and the stage slopes
\(k_i\). On a single step, write \(\tau=\tau_n\). As in the previous
lecture, \(\mathsf A=(a_{ij})\) denotes the RK coefficient matrix.

\subsection{Implicit Runge--Kutta methods and stability functions}

An \(s\)-stage RK method has the form
\begin{align}
 k_i&=f\left(t_n+c_i\tau,
       x_n+\tau\sum_{j=1}^s a_{ij}k_j\right),
       \qquad i=1,\ldots,s,
       \label{eq:rk-stages-sep-21}\\
 x_{n+1}&=x_n+\tau\sum_{i=1}^s b_i k_i.
       \label{eq:rk-update-sep-21}
\end{align}
The coefficients are arranged in the Butcher tableau
\[
 \begin{array}{c|c}
  c&\mathsf A\\ \hline
   &b^{\mathsf T}
 \end{array},
 \qquad
 c=\begin{bmatrix}c_1\\ \vdots\\ c_s\end{bmatrix},
 \quad
 b=\begin{bmatrix}b_1\\ \vdots\\ b_s\end{bmatrix}.
\]
For an implicit method, the stage equations generally form a coupled
algebraic system. They are nonlinear when \(f\) is nonlinear.

Apply the method to the scalar test equation
\[
 x'=\lambda x,\qquad \lambda\in\bb C,
 \qquad z=\tau\lambda.
\]
The stage equations become
\[
 k_i=\lambda x_n+z\sum_{j=1}^s a_{ij}k_j.
\]
With \(k=(k_1,\ldots,k_s)^{\mathsf T}\) and
\(\mathbf1=(1,\ldots,1)^{\mathsf T}\), this is
\[
 (I-z\mathsf A)k=\lambda x_n\mathbf1.
\]
Whenever \(I-z\mathsf A\) is invertible,
\[
 k=\lambda x_n(I-z\mathsf A)^{-1}\mathbf1.
\]
Substitution into the update gives
\begin{equation}\label{eq:rk-stability-function-sep-21}
 x_{n+1}=R(z)x_n,
 \qquad
 \boxed{R(z)=1+z b^{\mathsf T}(I-z\mathsf A)^{-1}\mathbf1.}
\end{equation}
Thus an implicit RK method generally has a rational stability
function. For an explicit method, \(\mathsf A^s=0\), so
\[
 (I-z\mathsf A)^{-1}
  =I+z\mathsf A+\cdots+z^{s-1}\mathsf A^{s-1},
\]
and \(R\) is a polynomial of degree at most \(s\).

Keeping the distinction from the previous lectures, the strict
absolute stability region and the non-strict stability domain are
\begin{align*}
 \mathcal S&=\{z\in\bb C: I-z\mathsf A\text{ is invertible},\ |R(z)|<1\},\\
 \mathcal D&=\{z\in\bb C: I-z\mathsf A\text{ is invertible},\ |R(z)|\leq1\}.
\end{align*}
For a fixed step, \(x_n=R(z)^n x_0\). Hence points in \(\mathcal S\)
give decay to zero, while points with \(|R(z)|=1\) give bounded
scalar iterates without decay. The invertibility condition ensures
that the stages themselves are uniquely defined.

\subsection{The collocation polynomial}

Consider a system \(x'=f(t,x)\), where the solution takes values in
\(\bb R^d\). Choose distinct nodes \(c_1,\ldots,c_s\in[0,1]\).
Let \(\mathcal P_m\) denote the
scalar polynomials of degree at most \(m\).

\begin{definition}[Collocation on one time step]
 Seek a vector polynomial \(u\in(\mathcal P_s)^d\) on
 \([t_n,t_{n+1}]\) satisfying
 \begin{align}
  u(t_n)&=x_n,\label{eq:collocation-initial-sep-21}\\
  u'(t_n+c_i\tau)
   &=f\bigl(t_n+c_i\tau,u(t_n+c_i\tau)\bigr),
   \qquad i=1,\ldots,s.
   \label{eq:collocation-equations-sep-21}
 \end{align}
 The next numerical value is
 \[
  x_{n+1}=u(t_{n+1}).
 \]
\end{definition}

The polynomial agrees with the initial value and satisfies the
differential equation at the chosen collocation points. It need not
satisfy the differential equation between those points. The initial
condition and the \(s\) collocation conditions give \((s+1)d\)
scalar equations for the \((s+1)d\) polynomial coefficients.
For smooth \(f\), the corresponding stage equations have a unique
local solution near \(\tau=0\); arbitrary large steps need not have
this property.

\subsection{Lagrange interpolation and the RK coefficients}

The Lagrange basis for the nodes is
\begin{equation}\label{eq:lagrange-basis-sep-21}
 \ell_j(\theta)
   =\prod_{\substack{m=1\\m\neq j}}^s
     \frac{\theta-c_m}{c_j-c_m},
 \qquad \ell_j(c_i)=\delta_{ij}.
\end{equation}
Each \(\ell_j\) has degree at most \(s-1\). Given scalar data
\(\varphi(c_j)\), the polynomial
\[
 P(\theta)=\sum_{j=1}^s\varphi(c_j)\ell_j(\theta)
\]
interpolates \(\varphi\) at all the nodes. In particular, if
\(\varphi\in\mathcal P_{s-1}\), then \(P=\varphi\) identically.

Define the collocation slopes by
\[
 k_j=u'(t_n+c_j\tau)
    =f\bigl(t_n+c_j\tau,u(t_n+c_j\tau)\bigr).
\]
Since \(u'\) has degree at most \(s-1\), its interpolation is exact:
\begin{equation}\label{eq:collocation-derivative-sep-21}
 u'(t_n+\theta\tau)=\sum_{j=1}^s k_j\ell_j(\theta).
\end{equation}
Integrating and using the initial condition gives
\begin{align}
 u(t_n+\theta\tau)-x_n
  &=\tau\int_0^\theta u'(t_n+\xi\tau)\,d\xi\notag\\
  &=\tau\sum_{j=1}^s
      \left(\int_0^\theta\ell_j(\xi)\,d\xi\right)k_j.
      \label{eq:collocation-polynomial-sep-21}
\end{align}
The factor \(\tau\) comes from the change of variables
\(t=t_n+\xi\tau\).

At \(\theta=c_i\), equation
\eqref{eq:collocation-polynomial-sep-21} yields
\[
 u(t_n+c_i\tau)=x_n+\tau\sum_{j=1}^s a_{ij}k_j,
 \qquad
 a_{ij}=\int_0^{c_i}\ell_j(\theta)\,d\theta.
\]
Inserting these stage values into
\eqref{eq:collocation-equations-sep-21} gives exactly the RK stage
equations \eqref{eq:rk-stages-sep-21}. At \(\theta=1\), we obtain
\[
 x_{n+1}=x_n+\tau\sum_{j=1}^s b_jk_j,
 \qquad
 b_j=\int_0^1\ell_j(\theta)\,d\theta.
\]
Consequently, the collocation coefficients are
\begin{equation}\label{eq:collocation-coefficients-sep-21}
 \boxed{
  a_{ij}=\int_0^{c_i}\ell_j(\theta)\,d\theta,
  \qquad
  b_j=\int_0^1\ell_j(\theta)\,d\theta.}
\end{equation}
Conversely, slopes satisfying this RK system define a collocation
polynomial through \eqref{eq:collocation-polynomial-sep-21}.
Thus the two formulations are equivalent.

Because \(\sum_j\ell_j(\theta)=1\), the coefficients automatically
satisfy
\[
 \sum_{j=1}^s a_{ij}=c_i,
 \qquad
 \sum_{j=1}^s b_j=1.
\]
These are the internal consistency relation
\(c=\mathsf A\mathbf1\) and the first-order condition
\(b^{\mathsf T}\mathbf1=1\).

Collocation generally produces implicit RK methods: the integrated
basis functions couple the unknown stages. The one-node choice
\(c_1=0\) is an exception, giving forward Euler. The choices
\(c_1=1\) and \((c_1,c_2)=(0,1)\) recover backward Euler and the
implicit trapezoidal method, respectively, as on September 18.

\subsection{The associated quadrature rule}

Integrating the Lagrange interpolant of a function \(\varphi\) gives
the interpolatory quadrature rule
\begin{equation}\label{eq:quadrature-rule-sep-21}
 \int_0^1\varphi(\theta)\,d\theta
   \approx\sum_{j=1}^s b_j\varphi(c_j),
 \qquad b_j=\int_0^1\ell_j(\theta)\,d\theta.
\end{equation}
The \(c_j\) are the quadrature nodes and the \(b_j\) are its weights.
Write the quadrature error as
\[
 E[\varphi]=\int_0^1\varphi(\theta)\,d\theta
              -\sum_{j=1}^s b_j\varphi(c_j).
\]
We use \(E\) here to distinguish the error functional from the
stability function \(R(z)\).

\begin{definition}[Degree of precision]
 A quadrature rule has \emph{degree of precision} \(p\) if
 \[
  E[\varphi]=0\quad\text{for every }\varphi\in\mathcal P_p,
 \]
 but \(E[\varphi]\neq0\) for at least one
 \(\varphi\in\mathcal P_{p+1}\).
\end{definition}

Every interpolatory rule with these \(s\) distinct nodes has
\(p\geq s-1\). Indeed, for \(\varphi\in\mathcal P_{s-1}\),
\[
 \varphi(\theta)=\sum_{j=1}^s\varphi(c_j)\ell_j(\theta),
\]
and integrating proves exactness. Equivalently, exactness through
degree \(p\) is expressed by the moment conditions
\begin{equation}\label{eq:quadrature-moments-sep-21}
 \sum_{j=1}^s b_jc_j^m=\frac1{m+1},
 \qquad m=0,\ldots,p.
\end{equation}
For \(m=0\), the left-hand side means \(\sum_j b_j\).

The largest possible degree of precision for \(s\) real nodes is
\(2s-1\). To see the upper bound, consider
\[
 q(\theta)=\prod_{j=1}^s(\theta-c_j)^2.
\]
This polynomial has degree \(2s\), vanishes at every node, and has
strictly positive integral on \([0,1]\). Therefore no such rule can
be exact for every polynomial of degree \(2s\).
Gaussian quadrature attains the bound \(p=2s-1\) with a suitable
choice of nodes. Its use in RK methods is the topic announced for
the next lecture.

\subsection{Quadrature precision and collocation order}

\begin{theorem}[Order of a collocation method]
 Suppose the interpolatory quadrature rule associated with the
 collocation nodes has degree of precision \(p\). Then the
 corresponding collocation RK method has order \(p+1\), assuming
 sufficient smoothness of \(f\) and the solution.
\end{theorem}

Here \(p\) denotes quadrature precision; the RK order is \(p+1\).
Thus a step started from the exact solution satisfies
\[
 x(t_{n+1})-u(t_{n+1})=O(\tau^{p+2}).
\]
Under the usual smoothness and Lipschitz assumptions, this leads to
global error \(O(h^{p+1})\) on a fixed finite interval, where
\(h=\max_n\tau_n\). The order statement concerns the values at step
endpoints; it does not assert the same error order at every point of
the collocation polynomial.

The role of quadrature is especially transparent when
\(f(t,x)=g(t)\) is independent of \(x\). Then
\[
 x_{n+1}=x_n+\tau\sum_{j=1}^s b_jg(t_n+c_j\tau),
\]
whereas the exact increment is
\[
 x(t_{n+1})-x(t_n)=\tau\int_0^1g(t_n+\theta\tau)\,d\theta.
\]
Taylor expansion of \(g\) and the moment conditions
\eqref{eq:quadrature-moments-sep-21} give a one-step error of
\(O(\tau^{p+2})\). The theorem extends this order conclusion to
general nonlinear collocation equations.

\begin{example}[One collocation node]
 With \(s=1\) and \(c_1=c\), the basis polynomial is
 \(\ell_1=1\), so the tableau is
 \[
  \begin{array}{c|c}c&c\\ \hline &1\end{array}.
 \]
 The quadrature rule is \(\int_0^1\varphi\approx\varphi(c)\).
 It is exact for linear polynomials precisely when \(c=1/2\).
 Thus the midpoint choice gives the second-order implicit midpoint
 method
 \[
  x_{n+1}=x_n+\tau f\left(t_n+\frac\tau2,
                              \frac{x_n+x_{n+1}}2\right).
 \]
 The endpoint choices \(c=0\) and \(c=1\) have precision zero
 and give the first-order forward and backward Euler methods.
\end{example}

\begin{example}[Endpoint collocation]
 For \(c_1=0\), \(c_2=1\), the weights are
 \(b_1=b_2=1/2\). The quadrature rule is exact for constants and
 linear polynomials, but
 \[
  \frac12(0^2+1^2)=\frac12\neq\frac13=\int_0^1\theta^2\,d\theta.
 \]
 Hence \(p=1\), and the associated trapezoidal collocation method
 has order two.
\end{example}

For Gaussian nodes, \(p=2s-1\), so the resulting \(s\)-stage Gauss
collocation method has order \(2s\). In particular, one Gauss node
gives order two and two Gauss nodes give order four. The placement
of the nodes, not just their number, determines the quadrature
precision and hence the collocation order.
